% intro.tex

%%%%%%%%%%%%%%%%%%%%
\begin{frame}{}
	\fig{width = 0.70\textwidth}{figs/verifying-history}
\end{frame}
%%%%%%%%%%%%%%%%%%%%

%%%%%%%%%%%%%%%%%%%%
\begin{frame}{}
  \begin{center}
    可序列化 (Serializability): 所有事务就像按照某个全序执行一样。

    \vspace{0.30cm}
    \fig{width = 0.70\textwidth}{figs/cobra-osdi2020}
  \end{center}
\end{frame}
%%%%%%%%%%%%%%%%%%%%

%%%%%%%%%%%%%%%%%%%%
\begin{frame}{}
  \begin{center}
    为了提高性能, 很多数据库仅提供快照隔离 (Snapshot Isolation)。

    \vspace{0.30cm}
    \fig{width = 0.80\textwidth}{figs/db-si}
  \end{center}
\end{frame}
%%%%%%%%%%%%%%%%%%%%

%%%%%%%%%%%%%%%%%%%%
\begin{frame}{Snapshot Isolation (SI)}
  \begin{center}
		\resizebox{0.50\textwidth}{!}{% isolation-si-write-skew-tikz.tex

\begin{tikzpicture}[
  node distance = 1.5cm and 1.5cm,
  txn/.style = {draw, inner sep = 3pt, align = center},
  comm/.style = {arrows = {Stealth-Stealth}, ultra thick, purple}]

  \node[txn] (t-alice)
    {$x_{1} \gets \readevent(\acctone)$\\[2pt]
     $x_{2} \gets \readevent(\accttwo)$\\[2pt]
     $\text{\bf if}\; x_{1} + x_{2} > 100$ \\[2pt]
     $\quad \red{x_{1} \gets x_{1} - 100}$ \\[2pt]
     $\quad \writeevent(\acctone, x_{1})$};

  \node[txn, right = of t-alice] (t-bob)
    {$x_{1} \gets \readevent(\acctone)$\\[2pt]
     $x_{2} \gets \readevent(\accttwo)$\\[2pt]
     $\text{\bf if}\; x_{1} + x_{2} > 100$ \\[2pt]
     $\quad \red{x_{2} \gets x_{2} - 100}$ \\[2pt]
     $\quad \writeevent(\accttwo, x_{2})$};

  \node[txn, right = of t-bob] (t-carol)
    {$x_{1} \gets \readevent(\acctone)$\\[2pt]
     $x_{2} \gets \readevent(\accttwo)$};

  \node[draw, fill = yellow!50, below = of t-bob, font = \large,
    inner sep = 8pt, minimum width = 260pt] (isolation)
    {{Snapshot Isolation (SI)}};

  \node[below = of isolation,
    label = {[label distance = 5pt]above : {\large $\acctone = \accttwo = 60$}}] (db)
    {\includegraphics[scale = 0.70]{figs/db-logo-dollar}};

  \draw[comm] (t-alice) -- (t-alice |- isolation.north);
  \draw[comm] (t-bob) -- (t-bob |- isolation.north);
  \draw[comm] (t-carol) -- (t-carol |- isolation.north);
  \draw[comm] (isolation) -- (isolation |- db.north);

  % \uncover<2->{
  \node[txn, below = 0.50cm of db, fill = green!30] (t-bob-execution)
    {$\readevent(\acctone, 60)$ \\[2pt]
     $\readevent(\accttwo, 60)$ \\[2pt]
     $\writeevent(\accttwo, -40)$};
  % }
  % \uncover<2->{
  \node[txn, left = of t-bob-execution] (t-alice-execution)
    {$\readevent(\acctone, 60)$ \\[2pt]
     $\readevent(\accttwo, 60)$ \\[2pt]
     $\writeevent(\acctone, -40)$};
  % }
  % \uncover<3->{
  \node[txn, right = of t-bob-execution, fill = red!30] (t-carol-execution)
    {$\readevent(\acctone, -40)$ \\[2pt]
     $\readevent(\accttwo, -40)$};

    \node[above = 1.50cm of t-carol-execution] (write-skew)
      {\green{\Large \textsc{Write Skew}}};
    \node[above = 0.80cm of t-carol-execution, scale = 3] (yes) {\yes};
  % }
\end{tikzpicture}}

    \vspace{0.20cm}
    \violet{Snapshot Read:} Each transaction reads data from a {\it snapshot} \\
      as of the time the transaction started.
  \end{center}
\end{frame}
%%%%%%%%%%%%%%%%%%%%

%%%%%%%%%%%%%%%%%%%%
\begin{frame}{Snapshot Isolation (SI)}
  \begin{center}
    \resizebox{0.48\textwidth}{!}{% isolation-si-lost-update-tikz.tex

\begin{tikzpicture}[
  node distance = 1.5cm and 1.5cm,
  txn/.style = {draw, inner sep = 3pt, align = center},
  comm/.style = {arrows = {Stealth-Stealth}, ultra thick, purple}]

  \node[txn] (t-alice)
    {$x \gets \readevent(\acct)$ \\[2pt]
     $x \gets x + 50$ \\[2pt]
     $\writeevent(\acct, x)$};

  \node[txn, right = of t-alice] (t-bob)
    {$x \gets \readevent(\acct)$ \\[2pt]
     $x \gets x + 25$ \\[2pt]
     $\writeevent(\acct, x)$};

  \node[txn, right = of t-bob] (t-carol)
    {$\readevent(\acct)$};

  \node[draw, fill = yellow!50, below = of t-bob, font = \large,
    inner sep = 8pt, minimum width = 260pt] (isolation)
    {{Snapshot Isolation (SI)}};

  \node[below = of isolation,
    label = {[label distance = 5pt]above : {\large $\acct = 0$}}] (db)
    {\includegraphics[scale = 0.70]{figs/db-logo-dollar}};

  \draw[comm] (t-alice) -- (t-alice |- isolation.north);
  \draw[comm] (t-bob) -- (t-bob |- isolation.north);
  \draw[comm] (t-carol) -- (t-carol |- isolation.north);
  \draw[comm] (isolation) -- (isolation |- db.north);

  % \uncover<2->{
  \node[txn, below = 0.50cm of db] (t-bob-execution)
    {$\readevent(\acct, 0)$ \\[2pt] $\writeevent(\acct, ?)$};
  \node[below = 0.50cm of db, scale = 3.00] (t-bob-abort) {\no};
  % }
  % \uncover<6->{
  % \node[txn, below = 0.50cm of db, fill = green!50] (t-bob-retry)
  %   {$\readevent(\acct, 0)$ \\[2pt] $\writeevent(\acct, 75)$};
  % }

  % \uncover<2->{
  \node[txn, left = of t-bob-execution] (t-alice-execution)
    {$\readevent(\acct, 0)$ \\[2pt] $\writeevent(\acct, 50)$};
  % }

  % \uncover<2->{
  \node[above right = 0.20cm and 1.00cm of t-bob-execution] (lost-update)
    {\red{\large \textsc{Lost Update}}};
  \node[above right = -0.60cm and 1.30cm of t-bob-execution, scale = 3] (no) {\no};
  % }
\end{tikzpicture}}
  \end{center}

  \vspace{-0.50cm}
  \violet{Snapshot Write:}
    Concurrent transactions {\it cannot} write to the same key.
    One of them must be aborted.
\end{frame}
%%%%%%%%%%%%%%%%%%%%

%%%%%%%%%%%%%%%%%%%%
\begin{frame}{}
  \begin{center}
    但是, 很多数据库\red{\bf 未能}正确实现快照隔离。

    \vspace{0.30cm}
    \fig{width = 0.85\textwidth}{figs/db-si-violations}
  \end{center}
\end{frame}
%%%%%%%%%%%%%%%%%%%%

%%%%%%%%%%%%%%%%%%%%
\begin{frame}{}
  \begin{center}
    高效的快照隔离检测算法

    \fig{width = 0.85\textwidth}{figs/polysi-vldb2023}
  \end{center}
\end{frame}
%%%%%%%%%%%%%%%%%%%%

%%%%%%%%%%%%%%%%%%%%
\begin{frame}{}
  \begin{center}
    \fig{width = 0.70\textwidth}{figs/checker-polysi}
  \end{center}
\end{frame}
%%%%%%%%%%%%%%%%%%%%